% ============================================================
% A27 — Deployment dan Operasional
% ============================================================
\flowmeta{A27}{Deployment dan Operasional}{Admin/Ops}

\begin{flowsection}{Tujuan}
Menggambarkan topologi rilis (deployment) dan operasional sistem: frontend,
backend, basis data, penyimpanan berkas, cron, serta backup.
\end{flowsection}

\begin{flowsection}{Prasyarat}
Akses server/CI dan kredensial lingkungan.
\end{flowsection}

\begin{flowsection}{Alur Rilis}
\begin{enumerate}
  \item Pengembangan pada branch fitur $\rightarrow$ merge ke \texttt{main} (development).
  \item Validasi (test backend, build frontend) $\rightarrow$ merge ke \texttt{master} (produksi).
  \item Deploy otomatis (CI/CD) membangun frontend \& menarik migrasi backend.
  \item Verifikasi smoke test setelah rilis.
\end{enumerate}
\end{flowsection}

\shot{gambar/admin/a27-1.png}{Topologi deployment dan operasional TUSKO}{fig:a27-1}

\begin{flowsection}{Operasional Rutin}
\begin{itemize}
  \item Pastikan cron \texttt{schedule:run} aktif.
  \item Pantau antrean (queue) dan log aplikasi.
  \item Lakukan backup basis data berkala.
  \item Perhatikan catatan internal: pembersihan data demo (T44) dan penghapusan mock/fallback (T45).
\end{itemize}
\end{flowsection}

\begin{flowsection}{Hasil yang Diharapkan}
Sistem rilis stabil dan operasional berjalan otomatis dengan pemantauan.
\end{flowsection}

\begin{flowsection}{Penanganan Error dan Kasus Khusus}
\begin{itemize}
  \item Build gagal $\rightarrow$ perbaiki sebelum merge.
  \item Migrasi gagal $\rightarrow$ tinjau log deploy dan ulangi.
\end{itemize}
\end{flowsection}

\begin{flowsection}{Kaitan Modul}
\texttt{backend/deploy/hostinger-deploy.sh}, \texttt{.github/workflows/}, \texttt{.env}.
\end{flowsection}

\docver{v1.0}{Sep 2026}
